\ifdefined\gkver\else\def\gkver{student}\fi
\documentclass[\gkver]{gaokaozhenti}
\begin{document}

\chapter{2019年上海}

\begin{enumerate}
\item
%% number: 1
%% paperName: 2019年上海等级考第 1 题 3 分
%% typeId: 1
%% type: 单选题
%% chapter: 直线运动
%% point: 加速度
%% method: 无
%% score: 3
%% degree: 850
%% duplicateId: 0
%% body:
以下运动中，物体加速度一定保持不变的是\xzanswer{C}
\fourchoices[answer=C]
{简谐运动}
{匀速圆周运动}
{竖直上抛运动}
{加速直线运动}
%% answer: C
%% memo:
\memoanswer{
本题原卷及材料中未提供详解，待补充。
}

\item
%% number: 2
%% paperName: 2019年上海等级考第 2 题 3 分
%% typeId: 1
%% type: 单选题
%% chapter: 原子和原子核
%% point: 原子核的组成
%% method: 无
%% score: 3
%% degree: 950
%% duplicateId: 0
%% body:
原子核内一定存在着\xzanswer{A}
\fourchoices[answer=A]
{质子}
{$\alpha$ 粒子}
{电子}
{光子}
%% answer: A
%% memo:
\memoanswer{
本题原卷及材料中未提供详解，待补充。
}

\item
%% number: 3
%% paperName: 2019年上海等级考第 3 题 3 分
%% typeId: 1
%% type: 单选题
%% chapter: 机械振动
%% point: 弹簧振子
%% method: 无
%% score: 3
%% degree: 750
%% duplicateId: 0
%% body:
弹簧振子做简谐运动，$t = 0$ 时振子经过平衡位置，机械能 $E$ 与时间 $t$ 的关系是\xzanswer{D}
\fourchoices[answer=D, ispicture=true, h=2.2cm]
{figs/03a.png}
{figs/03b.png}
{figs/03c.png}
{figs/03d.png}
%% answer: D
%% memo:
\memoanswer{
本题原卷及材料中未提供详解，待补充。
}

\item
%% number: 4
%% paperName: 2019年上海等级考第 4 题 3 分
%% typeId: 1
%% type: 单选题
%% chapter: 光学
%% point: 光的衍射
%% method: 无
%% score: 3
%% degree: 850
%% duplicateId: 0
%% body:
“泊松亮斑”是光的\xzanswer{B}
\onepicture{figs/04.png}
\fourchoices[answer=B]
{干涉现象，说明光具有波动性}
{衍射现象，说明光具有波动性}
{干涉现象，说明光具有粒子性}
{衍射现象，说明光具有粒子性}
%% answer: B
%% memo:
\memoanswer{
本题原卷及材料中未提供详解，待补充。
}

\item
%% number: 5
%% paperName: 2019年上海等级考第 5 题 3 分
%% typeId: 1
%% type: 单选题
%% chapter: 气体性质
%% point: 玻意耳定律
%% method: 无
%% score: 3
%% degree: 750
%% duplicateId: 0
%% body:
两容器中装有质量相等、温度相同的同种理想气体，若容积不同，则两容器中气体的\xzanswer{B}
\fourchoices[answer=B]
{分子平均动能相等，压强相等}
{分子平均动能相等，压强不等}
{分子平均动能不等，压强相等}
{分子平均动能不等，压强不等}
%% answer: B
%% memo:
\memoanswer{
本题原卷及材料中未提供详解，待补充。
}

\item
%% number: 6
%% paperName: 2019年上海等级考第 6 题 3 分
%% typeId: 1
%% type: 单选题
%% chapter: 曲线运动
%% point: 线速度角速度
%% method: 无
%% score: 3
%% degree: 750
%% duplicateId: 0
%% body:
在如图所示的传动装置中，内轮 a 和外轮 b 的转轴分别过各自圆心 O、O$'$ 且垂直纸面。轮 a 通过摩擦带动轮 b 做无相对滑动的转动，则两轮\xzanswer{C}
\onepicture{\includesvg{figs/06}}
\fourchoices[answer=C]
{转动方向相同，周期相等}
{转动方向相反，周期相等}
{转动方向相同，周期不等}
{转动方向相反，周期不等}
%% answer: C
%% memo:
\memoanswer{
本题原卷及材料中未提供详解，待补充。
}

\item
%% number: 7
%% paperName: 2019年上海等级考第 7 题 3 分
%% typeId: 1
%% type: 单选题
%% chapter: 力物体的平衡
%% point: 受力分析
%% method: 无
%% score: 3
%% degree: 750
%% duplicateId: 0
%% body:
如图，甲虫沿树枝由 A 缓慢上爬到 B 的过程中，树枝对甲虫的作用力\xzanswer{C}
\onepicture{\includesvg{figs/07}}
\fourchoices[answer=C]
{一直增大}
{一直减小}
{保持不变}
{先减小后增大}
%% answer: C
%% memo:
\memoanswer{
本题原卷及材料中未提供详解，待补充。
}

\item
%% number: 8
%% paperName: 2019年上海等级考第 8 题 3 分
%% typeId: 1
%% type: 单选题
%% chapter: 机械波
%% point: 干涉衍射
%% method: 无
%% score: 3
%% degree: 558
%% duplicateId: 0
%% body:
两波源 Ⅰ、Ⅱ 在水槽中形成的波形如图所示，其中实线为波峰，虚线为波谷，则\xzanswer{D}
\onepicture{\includesvg{figs/08}}
\fourchoices[answer=D]
{a 是振动始终加强的点}
{b 是振动始终加强的点}
{a 和 b 都是振动始终加强的点}
{a 和 b 都不是振动始终加强的点}
%% answer: D
%% memo:
\memoanswer{
本题原卷及材料中未提供详解，待补充。
}

\item
%% number: 9
%% paperName: 2019年上海等级考第 9 题 4 分
%% typeId: 1
%% type: 单选题
%% chapter: 气体性质
%% point: 气体图象
%% method: 无
%% score: 4
%% degree: 650
%% duplicateId: 0
%% body:
一定质量的理想气体，其状态经历 $a\to b\to c$ 的变化，$p-T$ 图线如图所示，在该过程中气体体积\xzanswer{D}
\onepicture{\includesvg{figs/09}}
\fourchoices[answer=D]
{先不变后增大}
{先增大后不变}
{先不变后减小}
{先减小后不变}
%% answer: D
%% memo:
\memoanswer{
本题原卷及材料中未提供详解，待补充。
}

\item
%% number: 10
%% paperName: 2019年上海等级考第 10 题 4 分
%% typeId: 1
%% type: 单选题
%% chapter: 电磁感应
%% point: 楞次定律
%% method: 无
%% score: 4
%% degree: 650
%% duplicateId: 0
%% body:
如图，直流电路与导体圆环 a、b 共面，电路由电源 E 和滑动变阻器 R 构成。在滑动变阻器的滑片向上滑动过程中，a、b 中的感应电流方向是\xzanswer{D}
\onepicture{\includesvg{figs/10}}
\fourchoices[answer=D]
{a 顺时针，b 顺时针}
{a 逆时针，b 逆时针}
{a 顺时针，b 逆时针}
{a 逆时针，b 顺时针}
%% answer: D
%% memo:
\memoanswer{
本题原卷及材料中未提供详解，待补充。
}

\item
%% number: 11
%% paperName: 2019年上海等级考第 11 题 4 分
%% typeId: 1
%% type: 单选题
%% chapter: 磁场
%% point: 左手定则
%% method: 无
%% score: 4
%% degree: 650
%% duplicateId: 0
%% body:
如图，在薄金属圆筒表面上通以与其轴线平行、分布均匀的恒定电流时，该圆筒的形变趋势为\xzanswer{C}
\onepicture{\includesvg{figs/11}}
\fourchoices[answer=C]
{沿轴线上下压缩}
{沿轴线上下拉伸}
{沿半径向内收缩}
{沿半径向外膨胀}
%% answer: C
%% memo:
\memoanswer{
本题原卷及材料中未提供详解，待补充。
}

\item
%% number: 12
%% paperName: 2019年上海等级考第 12 题 4 分
%% typeId: 1
%% type: 单选题
%% chapter: 直线运动
%% point: 自由落体运动
%% method: 逆向法
%% score: 4
%% degree: 550
%% duplicateId: 0
%% body:
将演员从高处跳下的视频倒序播放，可实现演员轻松跃上高处的特效。在观众看来，演员上升过程中在低处的高度变化\xzanswer{A}
\fourchoices[answer=A]
{比高处快；加速度向下}
{比高处慢；加速度向下}
{比高处快；加速度向上}
{比高处慢；加速度向上}
%% answer: A
%% memo:
\memoanswer{
本题原卷及材料中未提供详解，待补充。
}

\item
%% number: 13
%% paperName: 2019年上海等级考第 13 题 4 分
%% typeId: 3
%% type: 填空题
%% chapter: 运动和力
%% point: 牛顿第二定律
%% method: 无
%% score: 4
%% degree: 756
%% duplicateId: 0
%% body:
“神舟十号太空授课”的一项内容是测量航天员的“体重”。测量设备以恒定拉力拉动航天员，并同时测量两个物理量：设备提供的拉力和\tkanswer[2.5cm]{宇航员的加速度}；再根据\tkanswer[2.5cm]{牛顿第二定律}就能得出航天员的质量。
%% answer: 宇航员的加速度，牛顿第二定律
%% memo:
\memoanswer{
本题原卷及材料中未提供详解，待补充。
}

\item
%% number: 14
%% paperName: 2019年上海等级考第 14 题 4 分
%% typeId: 3
%% type: 填空题
%% chapter: 电场
%% point: 电场线
%% method: 无
%% score: 4
%% degree: 650
%% duplicateId: 0
%% body:
如图电场中，P、Q 两点的电场强度大小 $E_{P} > E_{Q}$，其理由是\tkanswer[4cm]{P 点所在位置处的电场线较 Q 点处的密}；电势 $\varphi_{P} > \varphi_{Q}$，其理由是\tkanswer[5.5cm]{P、Q 两点位于同一根电场线上，由于沿着电场线方向电势降低，所以 P 点电势高}。
\onepicture[align=right]{\includesvg{figs/14}}
%% answer: P 点所在位置处的电场线较 Q 点处的密；P、Q 两点位于同一根电场线上，由于沿着电场线方向电势降低，所以 P 点电势高
%% memo:
\memoanswer{
本题原卷及材料中未提供详解，待补充。
}

\item
%% number: 15
%% paperName: 2019年上海等级考第 15 题 4 分
%% typeId: 3
%% type: 填空题
%% chapter: 磁场
%% point: 磁场的叠加
%% method: 无
%% score: 4
%% degree: 550
%% duplicateId: 0
%% body:
已知环形电流在圆心处的磁感应强度大小与其半径成反比，通电直导线在其延长线上的磁感应强度为零。纸面内闭合导线框 a、b 由同心半圆与直线构成，通以大小相同的电流，如图所示。a、b 在其圆心 O$_{a}$、O$_{b}$ 处的磁感应强度的大小 $B_{a}$\tkanswer[1cm]{$<$}$B_{b}$（选填“$>$”“=”或“$<$”），O$_{a}$ 处的磁场方向\tkanswer[2.5cm]{垂直纸面向上}。
\onepicture[align=right]{\includesvg{figs/15}}
%% answer: $<$，垂直纸面向上
%% memo:
\memoanswer{
本题原卷及材料中未提供详解，待补充。
}

\item
%% number: 16
%% paperName: 2019年上海等级考第 16 题 4 分
%% typeId: 3
%% type: 填空题
%% chapter: 气体性质
%% point: 气体的状态参量
%% method: 无
%% score: 4
%% degree: 650
%% duplicateId: 0
%% body:
如图，硬质塑料瓶灌入水后盖紧瓶盖，在下端开一小孔后水开始流出，且水流逐渐减小，瓶中气体经历的过程可近似视为\tkanswer[1.2cm]{等温}过程；水流停止后，水面与孔之间的高度差为 $h$，瓶中气体的压强为\tkanswer[2.5cm]{$p_{0} - \rho gh$}。（已知大气压强为 $p_{0}$，水的密度为 $\rho$，重力加速度 $g$）
\onepicture[align=right]{\includesvg{figs/16}}
%% answer: 等温；$p_{0} - \rho gh$
%% memo:
\memoanswer{
本题原卷及材料中未提供详解，待补充。

【参考】课本上的第六章 D 压缩气体的应用一节的自主活动：在一个如图所示的扁平塑料饮料瓶靠近底部的四面，用针对称地刺四个小孔，灌入大半瓶水后，比较盖紧瓶盖和打开瓶盖两种情况下所见现象的区别。再从不同方向挤压瓶体，观察所发生的情况，并作出解释。

\onepicture[w=4cm]{figs/16_an.jpg}
}

\item
%% number: 17
%% paperName: 2019年上海等级考第 17 题 4 分
%% typeId: 3
%% type: 填空题
%% chapter: 电路
%% point: 动态分析
%% method: 无
%% score: 4
%% degree: 550
%% duplicateId: 0
%% body:
在如图电路中，电源电动势为 $E$，内阻为 $r$，$R_{0}$ 为定值电阻。当滑动变阻器 R 的滑片向右移动时，电流表 A 的示数如何变化？

分析思路为：①R 接入电路的阻值增大$\to$②回路总电阻增大$\to$③总电流 $I$ 减小$\to$④端电压 $U$ 增大$\to$⑤流过 $R_{0}$ 的电流 $I_{0}$ 增大$\to$⑥电流表中的电流 $I_{A}$ 减小

由③推得④的理由是：\tkanswer[2.5cm]{$U = E - Ir$}；由⑤推得⑥的理由是：\tkanswer[3cm]{$I_{A} = I - I_{0}$}。
\onepicture[align=right]{\includesvg{figs/17}}
%% answer: $U = E - Ir$，$I_{A} = I - I_{0}$
%% memo:
\memoanswer{
本题原卷及材料中未提供详解，待补充。
}

\item
%% number: 18
%% paperName: 2019年上海等级考第 18 题 10 分
%% typeId: 4
%% type: 实验题
%% chapter: 机械振动
%% point: 单摆测重力加速度
%% method: 无
%% score: 10
%% degree: 550
%% duplicateId: 0
%% body:
做“用单摆测定重力加速度”的实验，供选择的器材有：

A．$1.2\Um$ 长的细线；B．约 $2\Um$ 长的弹性绳；C．带孔的小铁球；D．带孔的软木球；E．光电门
\begin{enumerate}
	\item
	应选用的摆线是\tkanswer[1cm]{A}，摆球是\tkanswer[1cm]{C}（均选填器材前面的字母）。光电门应放在摆球摆动过程中的\tkanswer[1.5cm]{最低点}。
	\item
	单摆摆动过程中光电门接收到光信号强度 $I$ 随时间 $t$ 的变化如图所示，由此可得单摆周期为\tkanswer[1cm]{C}。

	\fourchoices[answer=C]
	{$t_{1}-0$}
	{$t_{2}-t_{1}$}
	{$t_{3}-t_{1}$}
	{$t_{4}-t_{1}$}
	\item
	某同学利用秒表做此实验后，与用光电门做实验的同学们交流，发现自己测得的重力加速度值明显偏大。该同学在实验中出现的操作失误可能是\tkanswer[2.5cm]{计数次数偏大}。
\end{enumerate}
\onepicture[align=right]{\includesvg{figs/18}}
%% answer: 
\jdanswer{
\begin{enumerate}
	\item
	A，C，最低点

	\item
	C

	\item
	计数次数偏大

\end{enumerate}
}
%% memo:
\memoanswer{
本题原卷及材料中未提供详解，待补充。
}

\item
%% number: 19
%% paperName: 2019年上海等级考第 19 题 14 分
%% typeId: 6
%% type: 计算题
%% chapter: 电磁感应
%% point: 电磁感应的电量
%% method: 无
%% score: 14
%% degree: 450
%% duplicateId: 0
%% body:
如图，半径为 $a$ 的金属圆环竖直固定，电阻忽略不计。空间存在与环面垂直的匀强磁场，磁感应强度为 $B$。一质量为 $m$ 的导体棒受到始终向上的拉力作用，以速率 $v$ 向下匀速运动，棒始终与 $y$ 轴垂直，扫过整个环面时与环接触良好。导体棒单位长度电阻值为 $r$，重力加速度为 $g$。求：
\begin{enumerate}
	\item
	导体棒中的感应电流 $I$ 及通过导体棒的总电荷量 $q$；
	\item
	导体棒位于 $y > 0$ 处拉力的功率 $P$。
\end{enumerate}
\onepicture[align=right]{\includesvg{figs/19}}
%% answer: 
\jdanswer{
\begin{enumerate}
	\item
	$I = \frac{Bv}{r}$，$q = \frac{2Ba}{r}$

	\item
	当 $0 < y < 2a$ 时：$P = \left(mg - \frac{2B^{2}v\sqrt{a^{2} - (a - y)^{2}}}{r}\right)v$；当 $y \geq 2a$ 时：$P = mgv$
\end{enumerate}
}
%% memo:
\memoanswer{
（1）$I = \frac{BLv}{R} = \frac{BLv}{Lr} = \frac{Bv}{r}$

$q = It = \frac{Bv}{r}\cdot\frac{2a}{v} = \frac{2Ba}{r}$

（2）当 $0 < y < 2a$ 时：$P = (mg - BIL)v = \left(mg - B\cdot\frac{Bv}{r}\cdot 2\sqrt{a^{2} - (a - y)^{2}}\right)v = \left(mg - \frac{2B^{2}v\sqrt{a^{2} - (a - y)^{2}}}{r}\right)v$

当 $y \geq 2a$ 时：$P = mgv$
}

\item
%% number: 20
%% paperName: 2019年上海等级考第 20 题 16 分
%% typeId: 6
%% type: 计算题
%% chapter: 功和能
%% point: 机械能守恒定律
%% method: 无
%% score: 16
%% degree: 450
%% duplicateId: 0
%% body:
如图，光滑轨道 abc 固定在竖直平面内，在 c 点与粗糙水平轨道 cd 相切，a、c 与轨道最低点 b 之间的高度差分别为 $H_{1}$ 和 $H_{2}$。一个质量为 $m$ 的小球 A 从 a 处由静止起沿轨道运动，在 b 处与质量也为 $m$ 的静止小滑块 B 相撞后静止，并将动能全部传递给 B。随后 B 沿轨道运动到 d 静止，从 c 运动到 d 的时间为 $t$。重力加速度为 $g$。
\begin{enumerate}
	\item
	求滑块与水平轨道 cd 间的动摩擦因数 $\mu$；
	\item
	可以用 A、B 在轨道 abc 上的运动来类比光电效应。分析说明小球 A、滑块 B 可以类比光电效应中的什么粒子，哪个物理量可以类比极限频率。
\end{enumerate}
\onepicture[align=right]{\includesvg{figs/20}}
%% answer: 
\jdanswer{
\begin{enumerate}
	\item
	$\mu = \frac{\sqrt{2g(H_{1} - H_{2})}}{gt}$

	\item
	光子，光电子，$H_{2}$。光电效应是指当光子入射到金属表面后，金属内表面附近的电子吸收了光子的全部能量后，脱离金属束缚的现象；在此过程中光子提供能量，电子吸收能量。因此（1）中的小球 A 可以类比为光子，滑块 B 可以类比成电子；A 与 B 的碰撞可以类比为电子吸收光子的能量；b、c 之间的高度差 $H_{2}$ 对应的能量差可类比极限频率对应的光子能量。因此 $H_{2}$ 可类比极限频率。
\end{enumerate}
}
%% memo:
\memoanswer{
（1）由 A 球运动到 b 处的机械能为：

$E = mgH_{1}$

然后将这部分能量传递给 B，设 B 球在 c 处的速度为 $v$，则由机械能守恒定律可知，

$mgH_{1} = mgH_{2} + \frac{1}{2}mv^{2}$

$v = \sqrt{2g(H_{1} - H_{2})}$

$v = at = \mu gt$

$\mu = \frac{\sqrt{2g(H_{1} - H_{2})}}{gt}$

（2）光子，光电子，$H_{2}$

光电效应是指当光子入射到金属表面后，金属内表面附近的电子吸收了光子的全部能量后，脱离金属束缚的现象；在此过程中光子提供能量，电子吸收能量。因此（1）中的小球 A 可以类比为光子，滑块 B 可以类比成电子；A 与 B 的碰撞可以类比为电子吸收光子的能量；b、c 之间的高度差 $H_{2}$ 对应的能量差可类比极限频率对应的光子能量。因此 $H_{2}$ 可类比极限频率。

由爱因斯坦光电方程 $E_{k} = h\nu - W_{0}$ 画出图像：

\onepicture[w=5.5cm]{figs/20_an1.png}

在此题中类似的有：$E_{k} = mgH_{1} - mgH_{2}$

\onepicture[w=5.5cm]{figs/20_an2.png}
}

\end{enumerate}

\end{document}
